\documentclass[../../../../hsm_tok_q.tex]{subfiles}

\begin{document}
    \textsf{図\ref{fig/hsm_tok_2022_2nd_01_07_01_q}}は，ある中学校の生徒31人について，%
    1か月に学校図書館を利用した回数を箱ひげ図に表したものである．

    \textsf{図\ref{fig/hsm_tok_2022_2nd_01_07_01_q}}から読み取れることとして正しいものを，%
    次の\textsf{ア}～\textsf{エ}のうちから選び，記号で答えよ．
    \\
    \begin{minipage}{\linewidth}
        \vspace{.5\baselineskip}
        {\centering
            \setlength{\abovecaptionskip}{2pt}
            \includegraphics%
                {\subfix{fig_hsm_tok_2022_2nd_01_07/fig_hsm_tok_2022_2nd_01_07_01_q.pdf}}
                \captionof{figure}{} \label{fig/hsm_tok_2022_2nd_01_07_01_q}
        }
    \end{minipage}

    \vspace{.3\baselineskip}
    \begin{itemize}[font=\textsf,topsep=2pt, leftmargin=1\zw,
        labelwidth=1.2\zw, labelsep=.4\zw, itemindent=1\zw, align=left]
        \item[ア] 学校図書館を9回以上利用した生徒は16人いる．
        \item[イ] 学校図書館を1回も利用しなかった生徒がいる．
        \item[ウ] 学校図書館を利用した回数の四分位範囲は19回である．
        \item[エ] 学校図書館を利用した回数が少ない方から数えて8番目の生徒は4回である．
    \end{itemize} 
\end{document}


